\documentclass[10pt,a4paper]{amsart}
\usepackage{amsmath,amssymb,amsthm,graphicx}
% Retained minimal preamble for the Stacks Project chapter more-algebra.tex (More on Algebra).
% The chapter file pinned-original-source.tex carries only \input{preamble} on its first line, so no
% preamble text lives inside the chapter itself.  The Stacks Project uses one single shared preamble for
% every chapter; the snapshot stages it identically as sources/stacks-chow/preamble.tex and
% sources/stacks-etale-cohomology/preamble.tex.  The block below is copied mechanically, line for line,
% from sources/stacks-chow/preamble.tex lines 14-18 (the Xy-pic setup) and 154-217 (the native theorem
% declarations together with the mathematical macros).  Nothing was added, removed or redefined; the
% wrapper class and the amsmath/amssymb/amsthm/graphicx base packages are supplied by the assembler.
% Chapter: more-algebra.tex -- More on Algebra.
% For commutative diagrams we use Xy-pic
\usepackage[all]{xy}

% We use 2cell for 2-commutative diagrams.
\xyoption{2cell}

% Theorem environments.
%
\theoremstyle{plain}
\newtheorem{theorem}[subsection]{Theorem}
\newtheorem{proposition}[subsection]{Proposition}
\newtheorem{lemma}[subsection]{Lemma}

\theoremstyle{definition}
\newtheorem{definition}[subsection]{Definition}
\newtheorem{example}[subsection]{Example}
\newtheorem{exercise}[subsection]{Exercise}
\newtheorem{situation}[subsection]{Situation}

\theoremstyle{remark}
\newtheorem{remark}[subsection]{Remark}
\newtheorem{remarks}[subsection]{Remarks}

\numberwithin{equation}{subsection}

% Macros
%
\def\lim{\mathop{\mathrm{lim}}\nolimits}
\def\colim{\mathop{\mathrm{colim}}\nolimits}
\def\Spec{\mathop{\mathrm{Spec}}}
\def\Hom{\mathop{\mathrm{Hom}}\nolimits}
\def\Ext{\mathop{\mathrm{Ext}}\nolimits}
\def\SheafHom{\mathop{\mathcal{H}\!\mathit{om}}\nolimits}
\def\SheafExt{\mathop{\mathcal{E}\!\mathit{xt}}\nolimits}
\def\Sch{\mathit{Sch}}
\def\Mor{\mathop{\mathrm{Mor}}\nolimits}
\def\Ob{\mathop{\mathrm{Ob}}\nolimits}
\def\Sh{\mathop{\mathit{Sh}}\nolimits}
\def\NL{\mathop{N\!L}\nolimits}
\def\CH{\mathop{\mathrm{CH}}\nolimits}
\def\proetale{{pro\text{-}\acute{e}tale}}
\def\etale{{\acute{e}tale}}
\def\QCoh{\mathit{QCoh}}
\def\Ker{\mathop{\mathrm{Ker}}}
\def\Im{\mathop{\mathrm{Im}}}
\def\Coker{\mathop{\mathrm{Coker}}}
\def\Coim{\mathop{\mathrm{Coim}}}

% Boxtimes
%
\DeclareMathSymbol{\boxtimes}{\mathbin}{AMSa}{"02}

%
% Macros for moduli stacks/spaces
%
\def\QCohstack{\mathcal{QC}\!\mathit{oh}}
\def\Cohstack{\mathcal{C}\!\mathit{oh}}
\def\Spacesstack{\mathcal{S}\!\mathit{paces}}
\def\Quotfunctor{\mathrm{Quot}}
\def\Hilbfunctor{\mathrm{Hilb}}
\def\Curvesstack{\mathcal{C}\!\mathit{urves}}
\def\Polarizedstack{\mathcal{P}\!\mathit{olarized}}
\def\Complexesstack{\mathcal{C}\!\mathit{omplexes}}
% \Pic is the operator that assigns to X its picard group, usage \Pic(X)
% \Picardstack_{X/B} denotes the Picard stack of X over B
% \Picardfunctor_{X/B} denotes the Picard functor of X over B
\def\Pic{\mathop{\mathrm{Pic}}\nolimits}
\def\Picardstack{\mathcal{P}\!\mathit{ic}}
\def\Picardfunctor{\mathrm{Pic}}
\def\Deformationcategory{\mathcal{D}\!\mathit{ef}}
\title{The Artin-Schreier type polynomial P(z) obtained from a primitive pth root of unity has coefficients in the ideal generated by w = 1 - zeta and satisfies an addition formula}
\author{The Stacks Project Authors}
\date{stacks-project:more-algebra.tex:local-snapshot-20261012; GNU FDL 1.2 or later}
\begin{document}
\maketitle
\setcounter{subsection}{1}
\noindent\textit{Editorial note.} All delivered mathematical text is the author's own, copied from the pinned Stacks Project chapter file more-algebra.tex; no statement or proof step was rewritten, reordered or removed.\par\medskip
\noindent\textit{Editorial note.} The chapter file pinned-original-source.tex carries only \textbackslash{}input\{preamble\} on its first line, so no preamble text lives inside the chapter. The retained file sources/stacks-more-algebra/retained-preamble.tex is a mechanical, line-for-line copy of the project's single shared preamble as staged in the snapshot; nothing was added, removed or redefined and no package was added.\par\medskip
\noindent\textit{Editorial note.} The source URL is the source-level Stacks Project URL, as required when more than one card is produced from one source; the per-claim tag is carried by claim\_key.\par\medskip
\noindent\textit{Editorial note.} The author writes 'The second part follows from a direct computation we omit'; that sentence is kept verbatim and marks the one step left to the reader. The proof also contains the author's typo 'the proof if the first part is done', which is reproduced unchanged rather than silently corrected.\par\medskip
\noindent\textit{Allowed standard prerequisites (editorial).}\begin{itemize}
\item The ring of integers Z[zeta] of a cyclotomic field and its unique prime above p. the reduction at proof lines 34992-34995 and the prime-ideal statement at proof lines 35003-35005. Supplies the unique prime (w) above p and hence the relation p = u w\textasciicircum{}\{p-1\} that drives the computation of the coefficients a\_i.
\item Divisibility of cyclotomic polynomials and the identity p = prod\_j (1 - zeta\textasciicircum{}j). the factorization and its evaluation at t = 1 at proof lines 34996-35002. Is the arithmetic input that identifies the valuation of p at w as p - 1.
\end{itemize}\medskip
\begin{lemma}
\label{lemma-prepare}
Let $A$ be a ring which contains a primitive $p$th root of unity $\zeta$.
Set $w = 1 - \zeta$. Then
$$
P(z) = \frac{(1 + wz)^p - 1}{w^p} =
z^p - z + \sum\nolimits_{0 < i < p} a_i z^i
$$
is an element of $A[z]$ and in fact $a_i \in (w)$. Moreover, we have
$$
P(z_1 + z_2 + w z_1 z_2) = P(z_1) + P(z_2) + w^p P(z_1) P(z_2)
$$
in the polynomial ring $A[z_1, z_2]$.
\end{lemma}

\begin{proof}
It suffices to prove this when
$$
A = \mathbf{Z}[\zeta] = \mathbf{Z}[x]/(x^{p - 1} + \ldots + x + 1)
$$
is the ring of integers of the cyclotomic field. The polynomial identity
$t^p - 1 = (t - 1)(t - \zeta) \ldots (t - \zeta^{p - 1})$
(which is proved by looking at the roots on both sides)
shows that
$t^{p - 1} + \ldots + t + 1 = (t - \zeta) \ldots (t - \zeta^{p - 1})$.
Substituting $t = 1$ we obtain
$p = (1 - \zeta)(1 - \zeta^2) \ldots (1 - \zeta^{p - 1})$.
The maximal ideal $(p, w) = (w)$ is the unique prime ideal of $A$
lying over $p$ (as fields of characteristic $p$ do not have nontrivial
$p$th roots of $1$). It follows that $p = u w^{p - 1}$ for some unit $u$.
This implies that
$$
a_i = \frac{1}{p} {p \choose i} u w^{i - 1}
$$
for $p > i > 1$ and $- 1 + a_1 = pw/w^p = u$. Since $P(-1) = 0$ we
see that $0 = (-1)^p - u$ modulo $(w)$. Hence $a_1 \in (w)$ and the
proof if the first part is done. The second part follows from a direct
computation we omit.
\end{proof}

\end{document}
